%%%%%%%%%%%%%%%%%%%%%%%%%%%%%%%%%%%%%%%%%%%%%%%%%%%%%%%%%%%%%%%%%%%%%%%%
% Plantilla TFG/TFM
% Universidad de Murcia. Facultad de Informática
% Realizado por: José Manuel Requena Plens
% Modificado: Pablo José Rocamora Zamora
% Contacto: pablojoserocamora@gmail.com
%%%%%%%%%%%%%%%%%%%%%%%%%%%%%%%%%%%%%%%%%%%%%%%%%%%%%%%%%%%%%%%%%%%%%%%%


\chapter*{Extended Abstract}
\addcontentsline{toc}{chapter}{Extended Abstract} %

Cryo-Electron Tomography (cryo-ET) is a structural biology imaging technique that enables the
three-dimensional reconstruction of biological specimens—such as individual proteins, protein
complexes, and intact cells—at near-native conditions and nanometric resolution. Unlike classical
structural methods such as X-ray crystallography or Nuclear Magnetic Resonance, cryo-ET does not
require sample crystallisation or chemical fixation; instead, specimens are rapidly vitrified in
liquid ethane to prevent ice crystal formation, preserving their native hydrated structure. During
acquisition, the electron beam traverses the sample at a series of tilt angles, typically ranging
from $-60°$ to $+60°$, and the resulting 2D projection images are computationally back-projected
and filtered to reconstruct a 3D density map known as a tomogram. Each voxel in the tomogram
encodes a local scattering density, which correlates with the distribution of atomic mass in the
corresponding region of the sample. Cryo-ET has become an indispensable tool in structural cell
biology for studying macromolecular machines, membrane architecture, and organelle organisation
in their cellular context.\\

The automated analysis of cryo-ET data has benefited enormously from deep learning methods over
the past decade. Convolutional neural networks and related architectures have been applied to tasks
such as tomogram denoising, organelle segmentation, sub-tomogram classification, and particle
picking. These supervised approaches, however, require large annotated training datasets: volumes
in which each voxel or region is labelled with its biological identity. Generating such ground-truth
labels experimentally is extraordinarily tedious and relies on expert manual curation, which
severely limits the scale at which labelled data can be produced. This data scarcity problem is
one of the primary bottlenecks in the application of deep learning to cryo-ET.\\

A well-established strategy to overcome this limitation is the generation of \textbf{synthetic
tomograms}: computationally produced volumes that simulate the physical image-formation process
of the electron microscope, including the projection geometry, the contrast transfer function,
missing-wedge artefacts arising from the limited tilt range, and realistic noise models.
When synthetic data is sufficiently faithful to experimental observations, networks trained on it
can generalise to real tomograms, reducing the need for manual annotation. The fidelity of
synthetic data depends critically on how realistically it reproduces the spatial distribution
and density of macromolecules within the simulation volume.\\

\textit{PolNet}~\cite{martinezSanchezPolnetGithub} is an open-source Python framework specifically designed for the generation of
synthetic cryo-ET datasets. It provides a modular simulation pipeline that models lipid membranes,
transmembrane protein complexes, and cytoplasmic macromolecular assemblies, ultimately producing
synthetic tomograms together with their ground-truth segmentation masks. The protein insertion
module of PolNet is based on \textit{SAWLC}~\cite{martinezSanchez2024simulating} (Self-Avoiding Worm-Like Chain), a stochastic
placement strategy in which each new protein is seeded near the surface of the most recently
inserted molecule. This surface-aware approach produces spatially clustered distributions that
resemble the crowded, heterogeneous organisation observed in real cells.Despite its biological motivation, SAWLC is computationally demanding. The sequential nature of the placement process, combined with high rejection rates in dense volumes, severely limits the algorithm's capability when modeling highly crowded environments, as the progressive fragmentation of the remaining free space prevents it from achieving a high degree of final occupancy. This throughput limitation
restricts the scale at which synthetic datasets can be produced and motivates the investigation of
more efficient protein insertion strategies.\\

The present Bachelor's Thesis addresses this bottleneck by adapting, implementing, and
benchmarking \textbf{Tetris}, a protein insertion algorithm that uses three-dimensional
FFT-based cross-correlation to determine, at each step, the optimal position for placing a
new macromolecule within the simulation volume. The algorithm is adapted from its original
formulation, extended to operate inside volumes that already contain biological structures
such as membranes, and evaluated in an independent repository with the explicit goal of
comparing its performance against SAWLC and assessing whether it constitutes a viable
replacement for the insertion module of PolNet.\\

The core idea behind Tetris is to reframe protein placement as an optimisation problem over
the space of all geometrically valid positions within a local search region. Given the current
state of the output volume and a new protein density map, the algorithm first extracts a local
bounding box around a randomly sampled seed coordinate. Within this bounding box, a binary
occupancy map is constructed by thresholding the existing density at a global isosurface value,
computed once as a fixed fraction of the minimum peak intensity across all protein types. A shell template is then generated by morphologically
dilating the binary protein mask and subtracting the protein itself: the resulting shell encodes
the ideal contact zone where a new protein should be placed to fit tightly against its neighbours
without overlapping. The three-dimensional cross-correlation between this shell template and the
binary occupancy map is computed using the Fast Fourier Transform (FFT), producing a score map
in which high values indicate positions where the new protein would make close contact with
existing structure. The position with the maximum valid score—subject to boundary constraints
and strict non-overlap verification—is selected, and the protein density is added to the output
volume at that location.\\

This formulation has two important consequences. First, each insertion finds the \textit{globally
best} valid position within the search region, rather than accepting the first non-overlapping
candidate discovered by random sampling. This leads to denser packing: proteins are placed in
a manner that minimises wasted interstitial space, achieving higher volumetric occupancy at the
end of the simulation. Second, the FFT-based score computation evaluates all candidate positions
simultaneously in $\mathcal{O}(N \log N)$ time, where $N$ is the number of voxels in the
bounding box, rather than requiring a separate overlap check for each candidate. This is
fundamentally more efficient than the iterative random-walk approach of SAWLC.\\

Beyond the core algorithm, three adaptations were required to integrate Tetris into the
PolNet pipeline. First, a \textbf{membrane pre-loading} step was introduced so that the
algorithm becomes aware of biological structures already present in the simulation volume,
such as lipid membranes. Before insertion begins, the voxels occupied by these
pre-existing structures are written into the Tetris output volume with a density value well
above the binarisation threshold, and a boolean \textit{allowed\_mask} is propagated to the
correlation step to penalise with $-10^{9}$ any candidate position that falls on a membrane
voxel. This guarantees that protein insertion is restricted to biologically permitted regions
without altering the internal logic of the placement routine. Second, an
\textbf{intelligent seed sampling} mechanism was added through a new \texttt{pick\_seed}
function that uniformly samples valid coordinates across the entire free volume rather than
starting from the geometric centre, which prevents the algorithm from leaving the tomogram
borders empty. Third, the algorithm was wrapped in a \textbf{cluster-growth loop}: after each
successful insertion the next seed is set to the position just occupied, encouraging dense
local packing; after each failure a new seed is drawn globally from \texttt{pick\_seed}, and
the protein type is considered exhausted once a configurable number of consecutive failures
(\texttt{TRIES\_CLUSTERING}) is reached.\\

A GPU-accelerated variant of Tetris was also developed using the \textit{CuPy} library,
which provides a NumPy-compatible API for array operations on NVIDIA graphics processing
units.The back-end is selected via a user-controlled boolean flag (\texttt{USE\_GPU}): when set to \texttt{True}, all heavy numerical operations (random rotation, gaussian smoothing,
morphological dilation, FFT-based cross-correlation and \texttt{argwhere}-based seed
sampling) are dispatched to \texttt{cupyx.scipy} on the GPU; otherwise the same code paths
fall back transparently to \texttt{scipy} on the CPU, without any change to the algorithmic
logic. The most significant performance gains come from two operations whose cost dominates
the CPU profile: the global \texttt{argwhere} over the entire volume inside \texttt{pick\_seed},
which is reduced from roughly two minutes to a fraction of a second by exploiting the massive
parallelism of the GPU, and the 3D FFT cross-correlation between the in-shell template and
the binarised local window, which is offloaded to \texttt{cupyx.scipy.signal.correlate}.
The 50-step inner validation loop that retrieves the maximum of the correlation map, checks
boundary constraints and verifies the absence of voxel overlap is preserved unchanged on the
GPU; this control flow remains sequential and forces a host-device synchronisation on each
\texttt{argmax} call, but its overall cost is comparable to that of pure host-side I/O
operations and does not negate the benefits obtained on the parallelisable stages. The net
effect is a global speedup of roughly $\approx \times 3$ over the CPU implementation on the
reference workload, with the two former bottlenecks effectively removed from the timing
profile.\\

To evaluate Tetris systematically, a benchmarking suite of three scripts was developed in the \texttt{benchmark/} directory, all sharing the same experimental configuration: eleven protein types at 10~Å resolution and a membrane volume derived from a real cryo-ET tomogram.

\texttt{profile\_test.py} instruments Tetris (CPU or GPU, controlled by \texttt{USE\_GPU}) with fine-grained timers and produces an occupancy-vs-time chart coloured by phase (seed selection, random rotation, template construction, FFT cross-correlation, overlap validation). Its purpose is to identify per-phase bottlenecks and guide GPU optimisation.

\texttt{profile\_comparison.py} runs Tetris GPU, Tetris CPU, and SAWLC sequentially on the full eleven-protein workload and produces a single occupancy-vs-time line chart together with a CSV log. The goal is to compare how quickly each algorithm fills the available volume.

\texttt{benchmark\_scientific\_comparison.py} sweeps a range of target occupancies from 2.5\% to 70\% in 2.5\% steps and runs both Tetris and SAWLC at each target, caching results in a JSON file to avoid redundant computation. It produces four comparative charts: packing fraction achieved vs target, execution time vs target, insertion throughput, and SAWLC failure count. The goal is to characterise the relative performance of both algorithms across the full density spectrum.

The \texttt{sawlc-cpu-gpu/} subdirectory complements the above with a log-based analysis. It stores pre-recorded execution logs for Tetris CPU, Tetris GPU, and SAWLC (one log per number of protein types, from 1 to 11) and contains \texttt{plot\_comparison.py}, which reads those logs and produces a four-panel figure (\texttt{resultados\_comparativa\_completa.png}): total occupancy reached, execution time, monomer population per protein type (stacked bars), and throughput (monomers per minute) — all as a function of the number of active protein types. This chart reveals which algorithm achieves the highest packing density and insertion rate as molecular diversity increases.\\

The preliminary results of this comparative study confirm the expected benefits of the Tetris
approach. Across all simulated scenarios, Tetris achieves higher volumetric occupancy than
SAWLC at the end of a full simulation run. The FFT-guided placement strategy packs proteins more
efficiently because it actively selects the contact position that maximises adjacency to existing
structure, whereas SAWLC's random walk may leave isolated free-volume pockets that are too
narrow to accommodate any remaining protein. This difference in packing quality is particularly
pronounced for smaller proteins, which can fill residual interstitial space that larger molecules
cannot access: the Tetris algorithm inserts substantially more monomers of small proteins than
SAWLC, contributing more to the final volumetric occupancy.\\       

In terms of execution time and throughput, the GPU variant of Tetris delivers the best
performance. The two dominant bottlenecks on CPU — the global \texttt{argwhere} inside
\texttt{pick\_seed} ($\approx 50\%$ of total time) and the FFT cross-correlation ($\approx 15\%$) — are
offloaded to the GPU via \texttt{CuPy}, reducing them to a negligible fraction of runtime.
The sequential overlap-checking loop is preserved unchanged on both backends. The net effect
is a speedup of $\times3$ over the CPU implementation. The CPU variant of Tetris is slower
than SAWLC in execution time, but reaches a significantly higher final occupancy because the
FFT-guided placement fills interstitial regions that SAWLC's random walk leaves inaccessible.\\

It is important to acknowledge that the two algorithms reflect different design philosophies.
SAWLC models proteins as a worm-like chain — each new molecule is seeded near the surface of
the previously inserted one, forming a polymer-like sequence that grows along a trajectory.
Tetris eliminates the polymer concept entirely: instead of a chain, it grows a compact cluster
that can start anywhere in the free volume and expands guided by the FFT correlation map,
adapting its shape to the available space at each step.\\

In conclusion, this work demonstrates that Tetris is a computationally efficient alternative
to SAWLC for protein insertion in synthetic cryo-ET tomogram generation. The algorithm achieves
higher packing densities and shorter runtimes through FFT-based global optimisation and
GPU-accelerated vectorised computation, positioning it as a strong candidate for replacing the
SAWLC module in PolNet pipelines where throughput and dataset scale are the primary concerns.
Future work will focus on incorporating biologically motivated spatial constraints into the
algorithm—for example, surface-proximate seeding similar to SAWLC—to improve the realism of
the generated distributions while retaining the computational advantages of the FFT-based
placement strategy. Further directions include extending the algorithm to handle elongated
filamentous assemblies, conducting a larger-scale evaluation across diverse tomogram geometries
and protein libraries, and exploring the use of Tetris as a pre-processing step that
initialises SAWLC simulations with a dense uniform packing before applying surface-aware
refinement.\\